\documentclass[10pt,a4paper,full]{physicsquiz}

\quiztitle{PHY 104 Vibrations and Wave Phenomena}
\quizcourse{100-Question Examination Practice Bank}
\quizauthor{Ot}
\quizinstitution{}
\quizversion{1.0}
\quizinstructions{Choose the best answer in each case. All questions carry one mark. A calculator may be used.}
\quizconstants{
  \(g=\SI{9.80}{m.s^{-2}}\),
  \(v_{\text{sound, air}}=\SI{343}{m.s^{-1}}\),
  \(c=\SI{3.00e8}{m.s^{-1}}\), and
  \(\rho_{\text{air}}=\SI{1.20}{kg.m^{-3}}\)
}

\newcommand{\bankq}[9]{%
  \begin{quizquestion}[
    id=#1,
    topic={#2},
    difficulty=#3,
    marks=1,
    correct=#4,
    tags={#5},
    outcome={#6}
  ]
    #7
    \begin{choiceoptions}
      #8
    \end{choiceoptions}
  \end{quizquestion}
  \begin{quizsolution}
    #9
  \end{quizsolution}
}

\begin{document}
\hypersetup{pageanchor=false}
\makequiztitle
\hypersetup{pageanchor=true}
\pagenumbering{arabic}

\begin{quizreferencecontent}
\section*{Curriculum Coverage}
\begin{multicols}{2}
\begin{itemize}[leftmargin=*]
  \item SHM and energy
  \item Damping, resonance, and transients
  \item Coupled SHM, quality factor, and power response
  \item Normal modes
  \item Wave properties and classification
  \item Interference, diffraction, dispersion, and polarization
  \item Wave equations in one, two, and three dimensions
  \item Echoes, beats, and the Doppler effect
  \item Propagation in gases, liquids, and solids
\end{itemize}
\end{multicols}
\end{quizreferencecontent}

\begin{quizbank}

% -----------------------------------------------------------------------------
% 1--12: SHM and energy (4 foundation, 5 applied, 3 challenge)
% -----------------------------------------------------------------------------

\bankq{shm-001}{SHM and Energy}{foundation}{B}{shm,definition}{Identify the defining acceleration law of SHM.}{%
Which relation defines ideal simple harmonic motion?}{%
\item \(a=+\omega^2x\)
\item \(a=-\omega^2x\)
\item \(a=-\omega x^2\)
\item \(a=\omega/x\)
\item \(a\) is constant}{%
In SHM the acceleration is proportional to displacement and directed toward equilibrium, hence \(a=-\omega^2x\).}

\bankq{shm-002}{SHM and Energy}{foundation}{D}{shm,kinematics}{Locate maximum speed in an SHM cycle.}{%
At which position is the speed of an ideal oscillator greatest?}{%
\item At \(x=A\)
\item At \(x=-A\)
\item At \(x=A/2\)
\item At \(x=0\)
\item At every position}{%
From \(v^2=\omega^2(A^2-x^2)\), the speed is largest when \(x=0\).}

\bankq{shm-003}{SHM and Energy}{foundation}{A}{shm,energy}{Recall the total energy of a spring oscillator.}{%
A spring oscillator has spring constant \(k\) and amplitude \(A\). Its total mechanical energy is}{%
\item \(\tfrac12kA^2\)
\item \(kA\)
\item \(\tfrac12mA^2\)
\item \(kA^2\)
\item \(\tfrac12kx^2\) at every position}{%
At a turning point \(v=0\) and \(|x|=A\), so all the energy is elastic: \(E=\tfrac12kA^2\).}

\bankq{shm-004}{SHM and Energy}{foundation}{C}{shm,period}{Relate period and angular frequency.}{%
The period of an oscillator with angular frequency \(\omega\) is}{%
\item \(1/\omega\)
\item \(\omega/2\pi\)
\item \(2\pi/\omega\)
\item \(2\pi\omega\)
\item \(\pi/\omega^2\)}{%
Since \(\omega=2\pi f\) and \(T=1/f\), \(T=2\pi/\omega\).}

\bankq{shm-005}{SHM and Energy}{applied}{C}{shm,spring,numerical}{Calculate a mass--spring period.}{%
A \(\SI{0.50}{kg}\) mass is attached to a spring of constant \(\SI{200}{N.m^{-1}}\). Its period is closest to}{%
\item \(\SI{0.157}{s}\)
\item \(\SI{0.222}{s}\)
\item \(\SI{0.314}{s}\)
\item \(\SI{0.628}{s}\)
\item \(\SI{1.57}{s}\)}{%
\(T=2\pi\sqrt{m/k}=2\pi\sqrt{0.50/200}=\SI{0.314}{s}\).}

\bankq{shm-006}{SHM and Energy}{applied}{E}{shm,energy,numerical}{Determine the potential-energy fraction in SHM.}{%
An oscillator has amplitude \(A\). At \(x=A/2\), what fraction of its total energy is elastic potential energy?}{%
\item \(1\)
\item \(3/4\)
\item \(1/2\)
\item \(1/3\)
\item \(1/4\)}{%
\(U/E=(\tfrac12kx^2)/(\tfrac12kA^2)=x^2/A^2=(1/2)^2=1/4\).}

\bankq{shm-007}{SHM and Energy}{applied}{B}{shm,kinematics,numerical}{Calculate speed from displacement in SHM.}{%
A particle has \(A=\SI{0.050}{m}\), \(\omega=\SI{10}{rad.s^{-1}}\), and \(x=\SI{0.030}{m}\). Its speed is}{%
\item \(\SI{0.20}{m.s^{-1}}\)
\item \(\SI{0.40}{m.s^{-1}}\)
\item \(\SI{0.50}{m.s^{-1}}\)
\item \(\SI{0.80}{m.s^{-1}}\)
\item \(\SI{1.00}{m.s^{-1}}\)}{%
\(v=\omega\sqrt{A^2-x^2}=10\sqrt{0.050^2-0.030^2}=\SI{0.40}{m.s^{-1}}\).}

\bankq{shm-008}{SHM and Energy}{applied}{D}{shm,pendulum,scaling}{Apply pendulum period scaling.}{%
The length of a small-angle pendulum is increased by a factor of four. Its new period is}{%
\item one-quarter of the original
\item one-half of the original
\item unchanged
\item twice the original
\item four times the original}{%
Because \(T=2\pi\sqrt{L/g}\), replacing \(L\) by \(4L\) multiplies \(T\) by \(2\).}

\bankq{shm-009}{SHM and Energy}{applied}{A}{shm,phase}{Interpret phase in an SHM expression.}{%
For \(x=A\cos(\omega t+\phi)\), a particle is at equilibrium and moving in the positive direction at \(t=0\). A suitable principal phase is}{%
\item \(-\pi/2\)
\item \(0\)
\item \(\pi/2\)
\item \(\pi\)
\item \(3\pi/2\) only}{%
At \(t=0\), \(x=A\cos\phi=0\), while \(v=-A\omega\sin\phi>0\). Thus \(\cos\phi=0\) and \(\sin\phi<0\), giving \(\phi=-\pi/2\).}

\bankq{shm-010}{SHM and Energy}{challenge}{C}{shm,energy,multistep}{Relate kinetic and potential energy to displacement.}{%
At what magnitude of displacement is the kinetic energy of an ideal spring oscillator three times its potential energy?}{%
\item \(A/4\)
\item \(A/\sqrt{3}\)
\item \(A/2\)
\item \(A/\sqrt{2}\)
\item \(\sqrt{3}A/2\)}{%
With \(K=E-U\) and \(K=3U\), \(E=4U\). Therefore \(\tfrac12kA^2=4(\tfrac12kx^2)\), so \(|x|=A/2\).}

\bankq{shm-011}{SHM and Energy}{challenge}{C}{shm,phase,time}{Determine a travel time between two SHM positions.}{%
A particle starts at \(x=A\) and moves toward equilibrium. How long does it first take to reach \(x=-A/2\)?}{%
\item \(T/6\)
\item \(T/4\)
\item \(T/3\)
\item \(5T/12\)
\item \(2T/3\)}{%
Using \(x=A\cos\omega t\), set \(\cos\omega t=-1/2\). The first solution is \(\omega t=2\pi/3\), hence \(t=(2\pi/3)/\omega=T/3\).}

\bankq{shm-012}{SHM and Energy}{challenge}{B}{shm,collision,energy}{Combine momentum conservation with SHM energy.}{%
A mass \(M\) moving at speed \(u\) sticks to an equal stationary mass attached to an unstretched spring of constant \(k\). What is the amplitude of the subsequent oscillation?}{%
\item \(u\sqrt{M/k}\)
\item \(u\sqrt{M/(2k)}\)
\item \(2u\sqrt{M/k}\)
\item \(u\sqrt{2M/k}\)
\item \(Mu/k\)}{%
Momentum gives the combined mass \(2M\) speed \(u/2\). Then \(\tfrac12(2M)(u/2)^2=\tfrac12kA^2\), so \(A=u\sqrt{M/(2k)}\).}

% -----------------------------------------------------------------------------
% 13--25: damping, resonance, and transients (4 foundation, 6 applied, 3 challenge)
% -----------------------------------------------------------------------------

\bankq{damping-013}{Damping, Resonance, and Transients}{foundation}{D}{damping,transient}{Distinguish transient and steady-state motion.}{%
In a damped driven oscillator, the transient motion is the part that}{%
\item always oscillates at the driving frequency
\item grows indefinitely
\item is independent of initial conditions
\item depends on initial conditions and decays with time
\item remains after an infinite time}{%
The homogeneous solution carries the initial-condition dependence and is multiplied by a decaying exponential, so it is the transient.}

\bankq{damping-014}{Damping, Resonance, and Transients}{foundation}{A}{damping,frequency,notation}{Distinguish driving and damped natural frequencies.}{%
In these notes, \(\omega_d\) in \(F_0\cos(\omega_dt)\) denotes the}{%
\item driving angular frequency
\item damped natural angular frequency
\item undamped natural angular frequency
\item decay constant
\item resonance width}{%
The argument of the applied force fixes its driving angular frequency, denoted here by \(\omega_d\).}

\bankq{damping-015}{Damping, Resonance, and Transients}{foundation}{C}{resonance,concept}{Recognize the resonance condition.}{%
For light damping, the displacement response is largest when the driving frequency is}{%
\item zero
\item much greater than the natural frequency
\item close to the natural frequency
\item equal to the damping constant
\item independent of the natural frequency}{%
Resonance occurs when the forcing rate is close to a natural oscillation rate of the system.}

\bankq{damping-016}{Damping, Resonance, and Transients}{foundation}{B}{damping,critical}{Identify critical damping.}{%
Critical damping is the damping that}{%
\item produces the longest possible oscillation
\item returns the system to equilibrium fastest without oscillation
\item makes the amplitude infinite
\item removes the restoring force
\item makes \(Q\) infinite}{%
Critical damping is the boundary between underdamped oscillation and overdamped non-oscillatory return, and gives the fastest non-oscillatory settling.}

\bankq{damping-017}{Damping, Resonance, and Transients}{applied}{E}{damping,decay,numerical}{Calculate a damping decay constant.}{%
The amplitude envelope is \(A=A_0e^{-\gamma t}\). If it falls to \(A_0/e\) in \(\SI{4.0}{s}\), \(\gamma\) is}{%
\item \(\SI{4.0}{s^{-1}}\)
\item \(\SI{2.0}{s^{-1}}\)
\item \(\SI{1.0}{s^{-1}}\)
\item \(\SI{0.50}{s^{-1}}\)
\item \(\SI{0.25}{s^{-1}}\)}{%
\(e^{-4\gamma}=e^{-1}\), so \(4\gamma=1\) and \(\gamma=\SI{0.25}{s^{-1}}\).}

\bankq{damping-018}{Damping, Resonance, and Transients}{applied}{C}{damping,frequency,numerical}{Calculate the damped natural frequency.}{%
For \(m=\SI{1.0}{kg}\), \(k=\SI{100}{N.m^{-1}}\), and \(b=\SI{4.0}{kg.s^{-1}}\), the damped natural angular frequency is}{%
\item \(\SI{2.00}{rad.s^{-1}}\)
\item \(\SI{8.00}{rad.s^{-1}}\)
\item \(\SI{9.80}{rad.s^{-1}}\)
\item \(\SI{10.0}{rad.s^{-1}}\)
\item \(\SI{12.0}{rad.s^{-1}}\)}{%
\(\gamma=b/(2m)=2\) and \(\omega_0=\sqrt{k/m}=10\). Hence \(\omega'=\sqrt{\omega_0^2-\gamma^2}=\sqrt{96}=\SI{9.80}{rad.s^{-1}}\).}

\bankq{damping-019}{Damping, Resonance, and Transients}{applied}{A}{resonance,amplitude,numerical}{Evaluate the driven amplitude at resonance.}{%
A system has \(F_0=\SI{5.0}{N}\), \(b=\SI{2.0}{kg.s^{-1}}\), and is driven at \(\omega_d=\SI{10}{rad.s^{-1}}\) where \(k-m\omega_d^2=0\). Its steady amplitude is}{%
\item \(\SI{0.25}{m}\)
\item \(\SI{0.50}{m}\)
\item \(\SI{1.0}{m}\)
\item \(\SI{2.0}{m}\)
\item \(\SI{25}{m}\)}{%
At this frequency \(X=F_0/(b\omega_d)=5/(2\times10)=\SI{0.25}{m}\).}

\bankq{damping-020}{Damping, Resonance, and Transients}{applied}{D}{resonance,phase}{Interpret phase at resonance.}{%
Near displacement resonance in a lightly damped oscillator, the displacement lags the driving force by approximately}{%
\item \(0\)
\item \(\pi/4\)
\item \(\pi\)
\item \(\pi/2\)
\item \(2\pi\)}{%
At resonance the reactive terms nearly cancel, so \(\tan\delta\) becomes large and \(\delta\approx\pi/2\).}

\bankq{damping-021}{Damping, Resonance, and Transients}{applied}{B}{transient,steady-state,frequency}{Identify frequencies in the complete driven response.}{%
An underdamped oscillator has \(\omega'=\SI{9.8}{rad.s^{-1}}\) and is driven at \(\omega_d=\SI{8.0}{rad.s^{-1}}\). Which statement is correct?}{%
\item Both parts oscillate permanently at \(\SI{9.8}{rad.s^{-1}}\)
\item The transient uses \(\omega'\), while the surviving steady state uses \(\omega_d\)
\item The transient uses \(\omega_d\), while the steady state uses \(\omega'\)
\item Both parts oscillate permanently at \(\SI{8.0}{rad.s^{-1}}\)
\item Neither frequency appears in the displacement}{%
The homogeneous transient oscillates at the damped natural frequency \(\omega'\); the particular steady response follows the driver at \(\omega_d\).}

\bankq{damping-022}{Damping, Resonance, and Transients}{applied}{D}{damping,critical,numerical}{Determine the critical damping coefficient.}{%
For \(m=\SI{2.0}{kg}\) and \(k=\SI{50}{N.m^{-1}}\), the critical damping coefficient is}{%
\item \(\SI{5}{kg.s^{-1}}\)
\item \(\SI{10}{kg.s^{-1}}\)
\item \(\SI{14.1}{kg.s^{-1}}\)
\item \(\SI{20}{kg.s^{-1}}\)
\item \(\SI{25}{kg.s^{-1}}\)}{%
\(b_c=2\sqrt{mk}=2\sqrt{2\times50}=\SI{20}{kg.s^{-1}}\).}

\bankq{damping-023}{Damping, Resonance, and Transients}{challenge}{B}{resonance,half-power,bandwidth}{Relate half-power frequencies to bandwidth and quality factor.}{%
A resonance curve has half-power frequencies \(\SI{98}{Hz}\) and \(\SI{102}{Hz}\), with resonance at \(\SI{100}{Hz}\). Its quality factor is}{%
\item \(4\)
\item \(25\)
\item \(50\)
\item \(100\)
\item \(200\)}{%
The bandwidth is \(\Delta f=102-98=\SI{4}{Hz}\). Thus \(Q=f_0/\Delta f=100/4=25\).}

\bankq{damping-024}{Damping, Resonance, and Transients}{challenge}{D}{resonance,power,derivation}{Locate maximum average power absorption.}{%
For \(m\ddot x+b\dot x+kx=F_0\cos\omega t\), at what angular frequency is the average absorbed power maximum?}{%
\item \(0\)
\item \(\sqrt{k/m-b^2/(2m^2)}\)
\item \(b/m\)
\item \(\sqrt{k/m}\)
\item It has no maximum}{%
The average power is \(\bar P=\tfrac12F_0^2b\omega^2/[(k-m\omega^2)^2+b^2\omega^2]\). Differentiation shows its maximum occurs at \(\omega=\omega_0=\sqrt{k/m}\).}

\bankq{damping-025}{Damping, Resonance, and Transients}{challenge}{C}{transient,resonance,superposition}{Analyze beating during the transient approach.}{%
A lightly damped oscillator with \(\omega'=\SI{10.0}{rad.s^{-1}}\) is driven at \(\omega_d=\SI{9.6}{rad.s^{-1}}\). Before the transient decays, the approximate angular beat rate is}{%
\item \(\SI{19.6}{rad.s^{-1}}\)
\item \(\SI{9.8}{rad.s^{-1}}\)
\item \(\SI{0.40}{rad.s^{-1}}\)
\item \(\SI{0.20}{rad.s^{-1}}\)
\item \(\SI{0.04}{rad.s^{-1}}\)}{%
The temporary response contains nearby frequencies \(\omega'\) and \(\omega_d\), so the beat angular rate is \(|\omega'-\omega_d|=|10.0-9.6|=\SI{0.40}{rad.s^{-1}}\).}

% -----------------------------------------------------------------------------
% 26--35: coupled SHM, Q values, and power response (3 foundation, 4 applied,
% 3 challenge)
% -----------------------------------------------------------------------------

\bankq{coupled-026}{Coupled SHM, Q, and Power Response}{foundation}{A}{coupling,concept}{Recognize the effect of coupling.}{%
The essential effect of coupling two oscillators is that it allows them to}{%
\item exchange energy
\item lose all restoring force
\item acquire identical displacements at all times
\item oscillate only at zero frequency
\item violate energy conservation}{%
The coupling force depends on relative motion and transfers energy between the oscillators.}

\bankq{coupled-027}{Coupled SHM, Q, and Power Response}{foundation}{D}{quality-factor,definition}{Interpret the quality factor.}{%
A large quality factor \(Q\) indicates}{%
\item very heavy damping and a broad resonance
\item no restoring force
\item a low natural frequency only
\item weak damping and a sharp resonance
\item that resonance cannot occur}{%
Large \(Q\) means little energy is lost per cycle, producing a narrow, high resonance peak.}

\bankq{coupled-028}{Coupled SHM, Q, and Power Response}{foundation}{B}{power,response-curve}{Identify half-power points.}{%
At a half-power frequency, the steady displacement amplitude is what fraction of its resonant value for a lightly damped oscillator?}{%
\item \(1/2\)
\item \(1/\sqrt{2}\)
\item \(1/4\)
\item \(\sqrt{2}\)
\item \(2\)}{%
Average power is proportional to amplitude squared near resonance. Half the maximum power therefore corresponds to amplitude \(X_{\max}/\sqrt{2}\).}

\bankq{coupled-029}{Coupled SHM, Q, and Power Response}{applied}{C}{quality-factor,numerical}{Calculate quality factor from oscillator parameters.}{%
An oscillator has \(m=\SI{0.50}{kg}\), \(\omega_0=\SI{20}{rad.s^{-1}}\), and \(b=\SI{0.40}{kg.s^{-1}}\). Its quality factor is}{%
\item \(5\)
\item \(10\)
\item \(25\)
\item \(40\)
\item \(50\)}{%
For viscous damping, \(Q=m\omega_0/b=(0.50)(20)/0.40=25\).}

\bankq{coupled-030}{Coupled SHM, Q, and Power Response}{applied}{E}{quality-factor,bandwidth}{Calculate resonance bandwidth.}{%
A resonator has \(f_0=\SI{600}{Hz}\) and \(Q=120\). Its half-power bandwidth is}{%
\item \(\SI{0.20}{Hz}\)
\item \(\SI{1.0}{Hz}\)
\item \(\SI{2.5}{Hz}\)
\item \(\SI{4.0}{Hz}\)
\item \(\SI{5.0}{Hz}\)}{%
\(Q=f_0/\Delta f\), so \(\Delta f=600/120=\SI{5.0}{Hz}\).}

\bankq{coupled-031}{Coupled SHM, Q, and Power Response}{applied}{A}{coupling,normal-frequency}{Calculate coupled-oscillator frequencies.}{%
Two identical masses have \(m=\SI{0.50}{kg}\), outer spring constant \(k=\SI{100}{N.m^{-1}}\), and coupling constant \(k_c=\SI{50}{N.m^{-1}}\). Their in-phase and out-of-phase angular frequencies are}{%
\item \(\SI{14.1}{rad.s^{-1}}\) and \(\SI{20.0}{rad.s^{-1}}\)
\item \(\SI{10.0}{rad.s^{-1}}\) and \(\SI{14.1}{rad.s^{-1}}\)
\item \(\SI{14.1}{rad.s^{-1}}\) and \(\SI{17.3}{rad.s^{-1}}\)
\item \(\SI{20.0}{rad.s^{-1}}\) and \(\SI{28.3}{rad.s^{-1}}\)
\item \(\SI{10.0}{rad.s^{-1}}\) and \(\SI{20.0}{rad.s^{-1}}\)}{%
\(\omega_+=\sqrt{k/m}=\sqrt{200}=\SI{14.1}{rad.s^{-1}}\), while \(\omega_-=\sqrt{(k+2k_c)/m}=\sqrt{400}=\SI{20.0}{rad.s^{-1}}\).}

\bankq{coupled-032}{Coupled SHM, Q, and Power Response}{applied}{D}{power,response,numerical}{Compare absorbed powers at one frequency.}{%
At the same driving frequency, one steady response has amplitude \(0.60\) times another. If the damping coefficient is unchanged, what is the ratio of their average absorbed powers?}{%
\item \(0.20\)
\item \(0.30\)
\item \(0.60\)
\item \(0.36\)
\item \(0.78\)}{%
At fixed \(b\) and \(\omega\), \(\bar P=\tfrac12b\omega^2X^2\). Therefore the power ratio is \((0.60)^2=0.36\).}

\bankq{coupled-033}{Coupled SHM, Q, and Power Response}{challenge}{B}{coupling,energy-transfer}{Calculate the first envelope-transfer time.}{%
Two coupled oscillators have normal angular frequencies \(\SI{14.14}{rad.s^{-1}}\) and \(\SI{20.00}{rad.s^{-1}}\). If only one is initially displaced, the first complete transfer of the oscillation envelope occurs after approximately}{%
\item \(\SI{0.268}{s}\)
\item \(\SI{0.536}{s}\)
\item \(\SI{1.07}{s}\)
\item \(\SI{2.14}{s}\)
\item \(\SI{5.86}{s}\)}{%
The transfer envelope first peaks at \(t=\pi/|\omega_2-\omega_1|=\pi/(20.00-14.14)=\SI{0.536}{s}\).}

\bankq{coupled-034}{Coupled SHM, Q, and Power Response}{challenge}{C}{quality-factor,energy-loss}{Relate Q to fractional energy loss.}{%
A lightly damped resonator has \(Q=50\). Approximately what fraction of its stored energy is lost per cycle?}{%
\item \(0.010\)
\item \(0.020\)
\item \(0.126\)
\item \(0.314\)
\item \(0.500\)}{%
For light damping, \(Q=2\pi(E/\Delta E)\). Thus \(\Delta E/E\approx2\pi/Q=2\pi/50=0.126\).}

\bankq{coupled-035}{Coupled SHM, Q, and Power Response}{challenge}{E}{power,response-curve,scaling}{Analyze how damping changes a resonance curve.}{%
The damping coefficient of a lightly damped driven oscillator is doubled while \(m\), \(k\), and \(F_0\) remain fixed. Which pair best describes the new resonant amplitude and half-power bandwidth?}{%
\item amplitude doubles; bandwidth halves
\item amplitude unchanged; bandwidth doubles
\item amplitude halves; bandwidth halves
\item amplitude becomes one-quarter; bandwidth doubles
\item amplitude approximately halves; bandwidth approximately doubles}{%
Near resonance \(X_{\max}\approx F_0/(b\omega_0)\), so doubling \(b\) halves the peak. Also \(\Delta\omega\approx b/m\), so the bandwidth doubles.}

% -----------------------------------------------------------------------------
% 36--45: normal modes (3 foundation, 4 applied, 3 challenge)
% -----------------------------------------------------------------------------

\bankq{modes-036}{Normal Modes}{foundation}{C}{normal-modes,definition}{Define a normal mode.}{%
In a normal mode, all coordinates oscillate}{%
\item with unrelated frequencies
\item with steadily changing phase differences
\item at one common frequency with fixed amplitude ratios and phases
\item only in the same direction
\item without exchanging any form of energy}{%
A normal mode is a collective pattern with one frequency and fixed relative motion among all coordinates.}

\bankq{modes-037}{Normal Modes}{foundation}{A}{normal-modes,in-phase}{Identify the in-phase eigenvector.}{%
For two identical coupled oscillators, the in-phase mode is proportional to}{%
\item \(\begin{bmatrix}1\\1\end{bmatrix}\)
\item \(\begin{bmatrix}1\\-1\end{bmatrix}\)
\item \(\begin{bmatrix}1\\0\end{bmatrix}\)
\item \(\begin{bmatrix}0\\1\end{bmatrix}\)
\item \(\begin{bmatrix}2\\-1\end{bmatrix}\)}{%
In phase means equal displacements with the same sign, represented by \([1,1]^T\).}

\bankq{modes-038}{Normal Modes}{foundation}{D}{normal-modes,eigenvalue}{Recall the normal-frequency condition.}{%
For \(M\ddot{\mathbf x}+K\mathbf x=0\), nonzero normal-mode amplitudes require}{%
\item \(\det(K+\omega M)=0\)
\item \(\det(M-\omega K)=0\)
\item \(\det(K-\omega M)=1\)
\item \(\det(K-\omega^2M)=0\)
\item \(K=M\)}{%
Substituting \(\mathbf x=\mathbf A\cos\omega t\) gives \((K-\omega^2M)\mathbf A=0\), which has nonzero solutions only when its determinant vanishes.}

\bankq{modes-039}{Normal Modes}{applied}{B}{normal-modes,matrix}{Find eigenfrequencies of a two-mode system.}{%
For unit masses and \(K=\begin{bmatrix}4&-2\\-2&4\end{bmatrix}\), the normal angular frequencies are}{%
\item \(1\) and \(2\)
\item \(\sqrt2\) and \(\sqrt6\)
\item \(2\) and \(\sqrt6\)
\item \(2\) and \(2\sqrt2\)
\item \(4\) and \(6\)}{%
\(\det(K-\omega^2I)=(4-\omega^2)^2-4=0\), so \(\omega^2=2,6\) and \(\omega=\sqrt2,\sqrt6\).}

\bankq{modes-040}{Normal Modes}{applied}{E}{normal-modes,decomposition}{Decompose an initial displacement into normal modes.}{%
Using mode vectors \([1,1]^T\) and \([1,-1]^T\), the displacement \([3,1]^T\) has respective coefficients}{%
\item \(3\) and \(1\)
\item \(1\) and \(3\)
\item \(4\) and \(2\)
\item \(1\) and \(1\)
\item \(2\) and \(1\)}{%
Solve \([3,1]^T=a[1,1]^T+b[1,-1]^T\). Then \(a+b=3\), \(a-b=1\), giving \(a=2\), \(b=1\).}

\bankq{modes-041}{Normal Modes}{applied}{C}{normal-modes,coupling,scaling}{Predict how coupling changes mode frequencies.}{%
For identical masses with fixed outer springs, increasing a positive coupling constant changes the two normal frequencies how?}{%
\item Both remain unchanged
\item Both increase equally
\item The in-phase frequency stays fixed while the out-of-phase frequency increases
\item The in-phase frequency decreases while the out-of-phase frequency stays fixed
\item Both become zero}{%
In phase, the coupling spring is unstrained, so \(\omega_+=\sqrt{k/m}\). Out of phase, it contributes restoring force, so \(\omega_-=\sqrt{(k+2k_c)/m}\) increases with \(k_c\).}

\bankq{modes-042}{Normal Modes}{applied}{A}{normal-modes,coordinates}{Use normal coordinates.}{%
If \(x_1=\SI{3}{cm}\) and \(x_2=-\SI{1}{cm}\), then \(q_+=x_1+x_2\) and \(q_-=x_1-x_2\) are}{%
\item \(\SI{2}{cm}\) and \(\SI{4}{cm}\)
\item \(\SI{4}{cm}\) and \(\SI{2}{cm}\)
\item \(\SI{2}{cm}\) and \(\SI{-4}{cm}\)
\item \(\SI{-2}{cm}\) and \(\SI{4}{cm}\)
\item \(\SI{3}{cm}\) and \(\SI{1}{cm}\)}{%
Direct substitution gives \(q_+=3+(-1)=\SI{2}{cm}\) and \(q_-=3-(-1)=\SI{4}{cm}\).}

\bankq{modes-043}{Normal Modes}{challenge}{D}{normal-modes,inverse-problem}{Infer spring constants from measured mode frequencies.}{%
A symmetric two-mass system has \(m=\SI{0.50}{kg}\), \(\omega_+=\SI{10}{rad.s^{-1}}\), and \(\omega_-=\SI{20}{rad.s^{-1}}\). The outer and coupling spring constants are}{%
\item \(k=\SI{25}{N.m^{-1}},\ k_c=\SI{50}{N.m^{-1}}\)
\item \(k=\SI{50}{N.m^{-1}},\ k_c=\SI{50}{N.m^{-1}}\)
\item \(k=\SI{100}{N.m^{-1}},\ k_c=\SI{25}{N.m^{-1}}\)
\item \(k=\SI{50}{N.m^{-1}},\ k_c=\SI{75}{N.m^{-1}}\)
\item \(k=\SI{200}{N.m^{-1}},\ k_c=\SI{75}{N.m^{-1}}\)}{%
From \(k=m\omega_+^2\), \(k=0.50(100)=\SI{50}{N.m^{-1}}\). Also \(k+2k_c=m\omega_-^2=0.50(400)=200\), so \(k_c=(200-50)/2=\SI{75}{N.m^{-1}}\).}

\bankq{modes-044}{Normal Modes}{challenge}{B}{normal-modes,initial-conditions}{Construct motion from two normal modes.}{%
For normal frequencies \(\omega_+\) and \(\omega_-\), the initial conditions \(x_1(0)=A\), \(x_2(0)=0\), and zero initial velocities give}{%
\item \(x_1=A\cos\omega_+t,\ x_2=0\)
\item \(x_1=\tfrac A2(\cos\omega_+t+\cos\omega_-t),\ x_2=\tfrac A2(\cos\omega_+t-\cos\omega_-t)\)
\item \(x_1=x_2=A\cos(\omega_++\omega_-)t\)
\item \(x_1=A\sin\omega_+t,\ x_2=A\sin\omega_-t\)
\item \(x_1=0,\ x_2=A\cos\omega_-t\)}{%
The initial vector \([A,0]^T\) is \(A[1,1]^T/2+A[1,-1]^T/2\). Evolving each mode at its own frequency gives option B.}

\bankq{modes-045}{Normal Modes}{challenge}{E}{normal-modes,generalized-eigenvalue}{Solve a generalized two-mass eigenvalue problem.}{%
For \(M=\begin{bmatrix}2&0\\0&1\end{bmatrix}\) and \(K=\begin{bmatrix}6&-2\\-2&4\end{bmatrix}\), the squared normal frequencies are}{%
\item \(1\) and \(2\)
\item \(2\) and \(3\)
\item \(1\) and \(5\)
\item \(3\) and \(4\)
\item \(2\) and \(5\)}{%
\(\det(K-\omega^2M)=(6-2s)(4-s)-4=2s^2-14s+20=2(s-2)(s-5)\), where \(s=\omega^2\). Thus \(\omega^2=2,5\), stated in option E.}

% -----------------------------------------------------------------------------
% 46--53: general wave properties and classification (3 foundation, 3 applied,
% 2 challenge)
% -----------------------------------------------------------------------------

\bankq{waves-046}{Wave Properties and Classification}{foundation}{B}{waves,classification}{Classify a transverse wave.}{%
In a transverse mechanical wave, particles of the medium oscillate}{%
\item parallel to the direction of propagation
\item perpendicular to the direction of propagation
\item only in circular paths
\item without acceleration
\item at the group velocity}{%
Transverse motion is defined by disturbance perpendicular to propagation.}

\bankq{waves-047}{Wave Properties and Classification}{foundation}{D}{waves,classification,sound}{Classify sound in a fluid.}{%
An ideal sound wave travelling through air is primarily}{%
\item transverse and polarized
\item electromagnetic
\item a surface wave
\item longitudinal
\item stationary}{%
Air supports alternating compressions and rarefactions, with molecular displacement parallel to propagation, so the wave is longitudinal.}

\bankq{waves-048}{Wave Properties and Classification}{foundation}{A}{waves,parameters}{Relate speed, frequency, and wavelength.}{%
The speed of a periodic wave is}{%
\item \(v=f\lambda\)
\item \(v=f/\lambda\)
\item \(v=\lambda/f\)
\item \(v=2\pi f/\lambda\)
\item \(v=\omega k\)}{%
One wavelength passes a point during one period, giving \(v=\lambda/T=f\lambda\).}

\bankq{waves-049}{Wave Properties and Classification}{applied}{C}{waves,wavefunction}{Extract wave parameters from a sinusoidal wave.}{%
For \(y=0.020\sin(4x-20t)\) in SI units, the amplitude and wave speed are}{%
\item \(\SI{4.0}{m}\) and \(\SI{20}{m.s^{-1}}\)
\item \(\SI{0.020}{m}\) and \(\SI{80}{m.s^{-1}}\)
\item \(\SI{0.020}{m}\) and \(\SI{5.0}{m.s^{-1}}\)
\item \(\SI{0.040}{m}\) and \(\SI{5.0}{m.s^{-1}}\)
\item \(\SI{0.020}{m}\) and \(\SI{0.20}{m.s^{-1}}\)}{%
The amplitude is \(0.020\) m. With \(k=4\) and \(\omega=20\), \(v=\omega/k=\SI{5.0}{m.s^{-1}}\).}

\bankq{waves-050}{Wave Properties and Classification}{applied}{E}{waves,phase}{Calculate separation from phase difference.}{%
Two points on a sinusoidal wave differ in phase by \(\pi/3\). If \(\lambda=\SI{0.60}{m}\), their smallest separation along the propagation direction is}{%
\item \(\SI{0.60}{m}\)
\item \(\SI{0.30}{m}\)
\item \(\SI{0.20}{m}\)
\item \(\SI{0.15}{m}\)
\item \(\SI{0.10}{m}\)}{%
\(\Delta\phi=2\pi\Delta x/\lambda\), so \(\Delta x=(\pi/3)(0.60)/(2\pi)=\SI{0.10}{m}\).}

\bankq{waves-051}{Wave Properties and Classification}{applied}{B}{waves,direction}{Determine propagation direction from phase.}{%
The wave \(y=A\cos(kx+\omega t)\) travels in the}{%
\item \(+x\) direction
\item \(-x\) direction
\item \(+y\) direction
\item \(-y\) direction
\item direction of increasing amplitude}{%
Constant phase requires \(kx+\omega t=\text{constant}\), so \(x\) decreases as \(t\) increases. The wave travels toward \(-x\).}

\bankq{waves-052}{Wave Properties and Classification}{challenge}{D}{waves,energy,scaling}{Apply wave-power scaling.}{%
For a sinusoidal wave on the same string, the amplitude is doubled and the frequency is tripled. By what factor does the average transmitted power change?}{%
\item \(6\)
\item \(9\)
\item \(12\)
\item \(36\)
\item \(81\)}{%
For fixed string properties, \(P_{\text{av}}\propto A^2\omega^2\). The factor is \(2^2\times3^2=36\).}

\bankq{waves-053}{Wave Properties and Classification}{challenge}{A}{waves,particle-motion}{Distinguish particle velocity from wave velocity.}{%
For \(y=A\sin(kx-\omega t)\), the maximum transverse particle speed divided by the wave speed is}{%
\item \(Ak\)
\item \(A/k\)
\item \(k/A\)
\item \(A\omega/k^2\)
\item \(1\) for every wave}{%
The maximum particle speed is \(A\omega\), while the wave speed is \(\omega/k\). Their ratio is \(A\omega/(\omega/k)=Ak\).}

% -----------------------------------------------------------------------------
% 54--70: interference, diffraction, dispersion, and polarization
% (5 foundation, 8 applied, 4 challenge)
% -----------------------------------------------------------------------------

\bankq{phenomena-054}{Interference, Diffraction, Dispersion, and Polarization}{foundation}{C}{superposition,principle}{State the superposition principle.}{%
When two linear waves overlap, the resultant displacement is}{%
\item the product of their displacements
\item always zero
\item the algebraic sum of their individual displacements
\item the larger displacement only
\item independent of phase}{%
For a linear medium, simultaneous disturbances add point by point.}

\bankq{phenomena-055}{Interference, Diffraction, Dispersion, and Polarization}{foundation}{A}{interference,constructive}{Recognize constructive interference.}{%
Two coherent waves of equal frequency arrive in phase. Their interference is}{%
\item constructive
\item destructive
\item necessarily unpolarized
\item dispersive
\item incoherent}{%
In-phase displacements reinforce, producing constructive interference.}

\bankq{phenomena-056}{Interference, Diffraction, Dispersion, and Polarization}{foundation}{E}{diffraction,single-slit}{Recall the exact single-slit minimum condition.}{%
For an ideal slit of width \(a\), diffraction minima occur at}{%
\item \(a\cos\theta=m\lambda\)
\item \(2a\sin\theta=m\lambda\)
\item \(a\sin\theta=(m+\tfrac12)\lambda\)
\item \(a\tan\theta=m\lambda\)
\item \(a\sin\theta=m\lambda\), \(m=1,2,3,\ldots\)}{%
Destructive cancellation across the slit gives the exact ideal condition \(a\sin\theta=m\lambda\) for nonzero integer \(m\).}

\bankq{phenomena-057}{Interference, Diffraction, Dispersion, and Polarization}{foundation}{B}{polarization,classification}{Identify which waves can be polarized.}{%
Polarization is a property that directly demonstrates that a wave is}{%
\item longitudinal
\item transverse
\item stationary
\item necessarily mechanical
\item nondispersive}{%
Polarization selects an oscillation direction perpendicular to propagation, which is possible for transverse waves.}

\bankq{phenomena-058}{Interference, Diffraction, Dispersion, and Polarization}{foundation}{D}{dispersion,concept}{Define dispersion.}{%
Dispersion occurs when}{%
\item amplitude is independent of time
\item all frequencies travel at one speed
\item a wave reflects from a boundary
\item wave speed depends on frequency or wavelength
\item two waves have equal phase}{%
Frequency-dependent propagation speed causes the components of a pulse to separate or spread.}

\bankq{phenomena-059}{Interference, Diffraction, Dispersion, and Polarization}{applied}{B}{interference,amplitude}{Calculate a resultant amplitude.}{%
Two waves of equal amplitude \(A\) and phase difference \(2\pi/3\) overlap. Their resultant amplitude is}{%
\item \(0\)
\item \(A\)
\item \(\sqrt2A\)
\item \(\sqrt3A\)
\item \(2A\)}{%
\(A_R=2A|\cos(\phi/2)|=2A\cos(\pi/3)=A\).}

\bankq{phenomena-060}{Interference, Diffraction, Dispersion, and Polarization}{applied}{C}{interference,path-difference}{Classify interference from path difference.}{%
Two in-phase coherent sources have path difference \(3\lambda/2\) at a point. The interference there is}{%
\item fully constructive
\item partially constructive only
\item fully destructive
\item incoherent
\item independent of wavelength}{%
For in-phase sources, destructive interference occurs at \((m+1/2)\lambda\). Here \(3\lambda/2\) corresponds to \(m=1\).}

\bankq{phenomena-061}{Interference, Diffraction, Dispersion, and Polarization}{applied}{D}{diffraction,single-slit,numerical}{Calculate a first-minimum angle.}{%
Light of wavelength \(\SI{600}{nm}\) passes through a slit of width \(\SI{0.20}{mm}\). The first minimum is closest to}{%
\item \(0.043^\circ\)
\item \(0.086^\circ\)
\item \(0.115^\circ\)
\item \(0.172^\circ\)
\item \(0.344^\circ\)}{%
For the first minimum, \(\sin\theta=\lambda/a=3.00\times10^{-3}\). Thus \(\theta\approx0.00300\) rad \(=0.172^\circ\).}

\bankq{phenomena-062}{Interference, Diffraction, Dispersion, and Polarization}{applied}{A}{diffraction,scaling}{Predict how slit width changes diffraction.}{%
If the width of a single slit is halved, the small-angle width of its central diffraction maximum approximately}{%
\item doubles
\item halves
\item quadruples
\item remains unchanged
\item becomes zero}{%
The first minima satisfy \(\theta_1\approx\lambda/a\), so the central width \(2\theta_1\) is inversely proportional to \(a\).}

\bankq{phenomena-063}{Interference, Diffraction, Dispersion, and Polarization}{applied}{E}{polarization,malus-law}{Apply Malus's law to unpolarized light.}{%
Unpolarized light of intensity \(I_0\) passes through an ideal polarizer and then an analyzer at \(60^\circ\) to it. The final intensity is}{%
\item \(I_0\)
\item \(I_0/2\)
\item \(I_0/4\)
\item \(I_0/6\)
\item \(I_0/8\)}{%
The first polarizer gives \(I_1=I_0/2\). Malus's law then gives \(I_2=I_1\cos^260^\circ=(I_0/2)(1/4)=I_0/8\).}

\bankq{phenomena-064}{Interference, Diffraction, Dispersion, and Polarization}{applied}{C}{dispersion,phase-velocity,group-velocity}{Calculate phase and group velocities.}{%
A medium obeys \(\omega=3k+0.50k^2\). At \(k=\SI{4.0}{m^{-1}}\), \((v_p,v_g)\) is}{%
\item \((3,4)\,\mathrm{m\,s^{-1}}\)
\item \((4,5)\,\mathrm{m\,s^{-1}}\)
\item \((5,7)\,\mathrm{m\,s^{-1}}\)
\item \((7,5)\,\mathrm{m\,s^{-1}}\)
\item \((7,7)\,\mathrm{m\,s^{-1}}\)}{%
At \(k=4\), \(\omega=3(4)+0.5(16)=20\). Hence \(v_p=\omega/k=5\), while \(v_g=d\omega/dk=3+k=7\), in \(\mathrm{m\,s^{-1}}\).}

\bankq{phenomena-065}{Interference, Diffraction, Dispersion, and Polarization}{applied}{B}{dispersion,inverse-derivative}{Use the reciprocal derivative for group velocity.}{%
If \(k=B\omega^2\), the group velocity is}{%
\item \(2B\omega\)
\item \(1/(2B\omega)\)
\item \(B\omega\)
\item \(\omega/(2B)\)
\item \(2\omega/B\)}{%
Since \(v_g=d\omega/dk=(dk/d\omega)^{-1}\) and \(dk/d\omega=2B\omega\), \(v_g=1/(2B\omega)\).}

\bankq{phenomena-066}{Interference, Diffraction, Dispersion, and Polarization}{applied}{D}{diffraction,resolution}{Apply the Rayleigh criterion.}{%
A circular aperture of diameter \(\SI{5.0}{cm}\) observes light of wavelength \(\SI{500}{nm}\). Its Rayleigh angular resolution is closest to}{%
\item \(1.0\times10^{-7}\,\mathrm{rad}\)
\item \(2.0\times10^{-6}\,\mathrm{rad}\)
\item \(5.0\times10^{-6}\,\mathrm{rad}\)
\item \(1.22\times10^{-5}\,\mathrm{rad}\)
\item \(6.10\times10^{-5}\,\mathrm{rad}\)}{%
\(\theta_R=1.22\lambda/D=1.22(500\times10^{-9})/0.050=1.22\times10^{-5}\,\mathrm{rad}\).}

\bankq{phenomena-067}{Interference, Diffraction, Dispersion, and Polarization}{challenge}{A}{dispersion,implicit-differentiation}{Differentiate an implicit dispersion relation.}{%
A wave satisfies \(F(\omega,k)=\omega^2-c^2k^2-\alpha k^4=0\). Its group velocity is}{%
\item \((c^2k+2\alpha k^3)/\omega\)
\item \((c^2k+\alpha k^3)/\omega\)
\item \(\omega/k\)
\item \(2\omega/(c^2k+2\alpha k^3)\)
\item \(c^2k/\omega^2\)}{%
Implicit differentiation gives \(2\omega\,d\omega/dk-2c^2k-4\alpha k^3=0\). Therefore \(v_g=d\omega/dk=(c^2k+2\alpha k^3)/\omega\).}

\bankq{phenomena-068}{Interference, Diffraction, Dispersion, and Polarization}{challenge}{C}{interference,diffraction,missing-orders}{Determine missing double-slit orders.}{%
A double slit has slit separation \(d=4a\), where \(a\) is each slit width. Which interference maxima are suppressed by the single-slit diffraction minima?}{%
\item every second order
\item every third order
\item orders \(4,8,12,\ldots\)
\item orders \(5,10,15,\ldots\)
\item no orders}{%
Interference maxima obey \(d\sin\theta=m\lambda\), while diffraction minima obey \(a\sin\theta=p\lambda\). Thus \(m=p(d/a)=4p\), so every fourth order is missing.}

\bankq{phenomena-069}{Interference, Diffraction, Dispersion, and Polarization}{challenge}{D}{polarization,multiple-polarizers}{Analyze three successive polarizers.}{%
Unpolarized light passes through ideal polarizers at \(0^\circ\), \(30^\circ\), and \(90^\circ\). The transmitted intensity is}{%
\item \(I_0/2\)
\item \(3I_0/16\)
\item \(I_0/8\)
\item \(3I_0/32\)
\item \(3I_0/64\)}{%
After the first, \(I_1=I_0/2\). Then \(I_2=I_1\cos^230^\circ=3I_0/8\). The final axis differs by \(60^\circ\), so \(I_3=(3I_0/8)\cos^260^\circ=3I_0/32\).}

\bankq{phenomena-070}{Interference, Diffraction, Dispersion, and Polarization}{challenge}{B}{diffraction,secondary-maximum}{Use the exact single-slit secondary-maximum condition.}{%
For a slit of width \(\SI{0.20}{mm}\) and wavelength \(\SI{600}{nm}\), use the first root \(\beta=4.493\) of \(\tan\beta=\beta\). The first secondary maximum is closest to}{%
\item \(0.172^\circ\)
\item \(0.246^\circ\)
\item \(0.258^\circ\)
\item \(0.344^\circ\)
\item \(0.492^\circ\)}{%
From \(\beta=\pi a\sin\theta/\lambda\), \(\sin\theta=4.493(600\times10^{-9})/[\pi(0.20\times10^{-3})]=4.29\times10^{-3}\). Thus \(\theta\approx0.246^\circ\).}

% -----------------------------------------------------------------------------
% 71--80: wave equations in 1D, 2D, and 3D (3 foundation, 4 applied,
% 3 challenge)
% -----------------------------------------------------------------------------

\bankq{waveeq-071}{Wave Equations in 1D, 2D, and 3D}{foundation}{A}{wave-equation,one-dimensional}{Recognize the one-dimensional wave equation.}{%
The standard one-dimensional wave equation is}{%
\item \(u_{tt}=v^2u_{xx}\)
\item \(u_t=vu_x\)
\item \(u_{xx}=v^2u\)
\item \(u_{tt}=vu_x\)
\item \(u_t=v^2u_{xx}\)}{%
The second time derivative is proportional to the second spatial derivative: \(u_{tt}=v^2u_{xx}\).}

\bankq{waveeq-072}{Wave Equations in 1D, 2D, and 3D}{foundation}{C}{wave-equation,two-dimensional}{Identify the two-dimensional Laplacian.}{%
For a membrane displacement \(u(x,y,t)\), the spatial term in the wave equation is}{%
\item \(u_x+u_y\)
\item \(u_{xy}\)
\item \(u_{xx}+u_{yy}\)
\item \(u_{xx}-u_{yy}\)
\item \(u_xu_y\)}{%
The two-dimensional Laplacian is \(\nabla^2u=u_{xx}+u_{yy}\).}

\bankq{waveeq-073}{Wave Equations in 1D, 2D, and 3D}{foundation}{D}{wave-equation,boundary-condition}{Interpret a fixed-end boundary condition.}{%
A string fixed at \(x=0\) satisfies which boundary condition?}{%
\item \(u_x(0,t)=1\)
\item \(u_t(0,t)=0\) only
\item \(u_{xx}(0,t)=0\)
\item \(u(0,t)=0\)
\item \(u(0,t)=A\)}{%
A fixed support cannot move, so its displacement is zero for all time.}

\bankq{waveeq-074}{Wave Equations in 1D, 2D, and 3D}{applied}{B}{wave-equation,plane-wave}{Test a sinusoidal solution.}{%
The function \(u=A\cos(kx-\omega t)\) satisfies \(u_{tt}=v^2u_{xx}\) provided}{%
\item \(\omega=k/v\)
\item \(\omega=vk\)
\item \(\omega=v/k\)
\item \(\omega^2=vk\)
\item \(k=v\omega\)}{%
Substitution gives \(-\omega^2u=v^2(-k^2u)\), hence \(\omega^2=v^2k^2\) and \(\omega=vk\) for positive frequencies.}

\bankq{waveeq-075}{Wave Equations in 1D, 2D, and 3D}{applied}{E}{wave-equation,general-solution}{Recognize the d'Alembert solution.}{%
The general solution of the ideal one-dimensional wave equation can be written}{%
\item \(u=f(x)g(t)\) only
\item \(u=f(x-vt)g(x+vt)\)
\item \(u=f(vx-t)+g(vx+t)\)
\item \(u=f(x-vt)-f(x-vt)\)
\item \(u=f(x-vt)+g(x+vt)\)}{%
The two arbitrary functions represent right- and left-moving waves, respectively: \(f(x-vt)+g(x+vt)\).}

\bankq{waveeq-076}{Wave Equations in 1D, 2D, and 3D}{applied}{C}{wave-equation,speed}{Read wave speed from a PDE.}{%
A field obeys \(u_{tt}=144u_{xx}\) in SI units. Its wave speed is}{%
\item \(\SI{6}{m.s^{-1}}\)
\item \(\SI{10}{m.s^{-1}}\)
\item \(\SI{12}{m.s^{-1}}\)
\item \(\SI{72}{m.s^{-1}}\)
\item \(\SI{144}{m.s^{-1}}\)}{%
Comparing with \(u_{tt}=v^2u_{xx}\) gives \(v^2=144\), so \(v=\SI{12}{m.s^{-1}}\).}

\bankq{waveeq-077}{Wave Equations in 1D, 2D, and 3D}{applied}{A}{wave-equation,three-dimensional}{Interpret a three-dimensional plane wave.}{%
For \(u(\mathbf r,t)=A\cos(\mathbf k\cdot\mathbf r-\omega t)\), the direction of propagation is}{%
\item parallel to \(\mathbf k\)
\item perpendicular to \(\mathbf k\)
\item parallel to the amplitude \(A\)
\item always along \(+z\)
\item undefined}{%
Surfaces of constant phase move in the direction of increasing \(\mathbf k\cdot\mathbf r\), so propagation is parallel to \(\mathbf k\).}

\bankq{waveeq-078}{Wave Equations in 1D, 2D, and 3D}{challenge}{D}{wave-equation,string,derivation}{Connect string properties to the wave equation.}{%
A string under tension \(\tau\) and linear density \(\mu\) obeys \(u_{tt}=(\tau/\mu)u_{xx}\). If \(\tau=\SI{90}{N}\) and \(\mu=\SI{0.010}{kg.m^{-1}}\), its wave speed is}{%
\item \(\SI{9.0}{m.s^{-1}}\)
\item \(\SI{30}{m.s^{-1}}\)
\item \(\SI{90}{m.s^{-1}}\)
\item \(\SI{94.9}{m.s^{-1}}\)
\item \(\SI{9000}{m.s^{-1}}\)}{%
The coefficient equals \(v^2\), so \(v=\sqrt{\tau/\mu}=\sqrt{90/0.010}=\sqrt{9000}=\SI{94.9}{m.s^{-1}}\).}

\bankq{waveeq-079}{Wave Equations in 1D, 2D, and 3D}{challenge}{B}{wave-equation,membrane,normal-mode}{Calculate a rectangular-membrane mode frequency.}{%
A rectangular membrane of sides \(a=\SI{2.0}{m}\), \(b=\SI{1.0}{m}\) supports waves at \(v=\SI{20}{m.s^{-1}}\). The \((m,n)=(1,1)\) mode frequency is}{%
\item \(\SI{5.00}{Hz}\)
\item \(\SI{11.2}{Hz}\)
\item \(\SI{15.0}{Hz}\)
\item \(\SI{22.4}{Hz}\)
\item \(\SI{44.7}{Hz}\)}{%
For fixed edges, \(f_{mn}=\tfrac v2\sqrt{(m/a)^2+(n/b)^2}\). Thus \(f_{11}=10\sqrt{(1/2)^2+1^2}=\SI{11.2}{Hz}\).}

\bankq{waveeq-080}{Wave Equations in 1D, 2D, and 3D}{challenge}{C}{wave-equation,three-dimensional,dispersion}{Verify a three-dimensional plane-wave relation.}{%
A plane wave has \(\mathbf k=(2,3,6)\,\mathrm{m^{-1}}\) in a nondispersive medium with \(v=\SI{5.0}{m.s^{-1}}\). Its angular frequency is}{%
\item \(\SI{7}{rad.s^{-1}}\)
\item \(\SI{25}{rad.s^{-1}}\)
\item \(\SI{35}{rad.s^{-1}}\)
\item \(\SI{55}{rad.s^{-1}}\)
\item \(\SI{70}{rad.s^{-1}}\)}{%
\(|\mathbf k|=\sqrt{2^2+3^2+6^2}=7\,\mathrm{m^{-1}}\). For \(\omega=v|\mathbf k|\), \(\omega=5(7)=\SI{35}{rad.s^{-1}}\).}

% -----------------------------------------------------------------------------
% 81--90: acoustics, echoes, beats, and Doppler effect
% (3 foundation, 5 applied, 2 challenge)
% -----------------------------------------------------------------------------

\bankq{acoustics-081}{Acoustics: Echoes, Beats, and Doppler Effect}{foundation}{B}{acoustics,echo}{Interpret echo travel time.}{%
If an echo returns \(t\) seconds after a sound is emitted in a medium of speed \(v\), the reflector distance is}{%
\item \(vt\)
\item \(vt/2\)
\item \(2vt\)
\item \(v/(2t)\)
\item \(t/(2v)\)}{%
The sound travels to the reflector and back, a total distance \(2d=vt\). Therefore \(d=vt/2\).}

\bankq{acoustics-082}{Acoustics: Echoes, Beats, and Doppler Effect}{foundation}{D}{acoustics,beats}{Relate beat frequency to source frequencies.}{%
Two tones of frequencies \(f_1\) and \(f_2\) produce beats at frequency}{%
\item \(f_1+f_2\)
\item \((f_1+f_2)/2\)
\item \(2|f_1-f_2|\)
\item \(|f_1-f_2|\)
\item \(f_1f_2\)}{%
The slow amplitude modulation repeats at the absolute frequency difference: \(f_b=|f_1-f_2|\).}

\bankq{acoustics-083}{Acoustics: Echoes, Beats, and Doppler Effect}{foundation}{A}{acoustics,doppler}{Predict the sign of a Doppler shift.}{%
When a sound source approaches a stationary observer, the observed frequency is generally}{%
\item greater than the emitted frequency
\item smaller than the emitted frequency
\item always zero
\item unchanged for all source speeds
\item equal to the beat frequency}{%
An approaching source compresses the wavefront spacing, reducing wavelength and increasing the observed frequency.}

\bankq{acoustics-084}{Acoustics: Echoes, Beats, and Doppler Effect}{applied}{C}{acoustics,echo,numerical}{Calculate distance from an echo delay.}{%
An echo is heard \(\SI{0.80}{s}\) after a shout. Taking the sound speed as \(\SI{343}{m.s^{-1}}\), the reflecting wall is approximately}{%
\item \(\SI{68.6}{m}\) away
\item \(\SI{109}{m}\) away
\item \(\SI{137}{m}\) away
\item \(\SI{274}{m}\) away
\item \(\SI{429}{m}\) away}{%
\(d=vt/2=343(0.80)/2=\SI{137.2}{m}\), approximately \(\SI{137}{m}\).}

\bankq{acoustics-085}{Acoustics: Echoes, Beats, and Doppler Effect}{applied}{D}{acoustics,beats,tuning}{Infer an unknown tone from beat information.}{%
A tuning fork of frequency \(\SI{440}{Hz}\) produces \(\SI{6}{beats.s^{-1}}\) with a string. Tightening the string increases the beat rate. The original string frequency was}{%
\item \(\SI{428}{Hz}\)
\item \(\SI{434}{Hz}\)
\item \(\SI{440}{Hz}\)
\item \(\SI{446}{Hz}\)
\item \(\SI{452}{Hz}\)}{%
The original possibilities are \(440\pm6\), namely \(434\) or \(446\) Hz. Tightening raises the string frequency; if it was already \(446\) Hz, the difference would increase. Thus the original was \(\SI{446}{Hz}\).}

\bankq{acoustics-086}{Acoustics: Echoes, Beats, and Doppler Effect}{applied}{B}{acoustics,doppler,numerical}{Calculate a Doppler shift for a moving source.}{%
A source emitting \(\SI{500}{Hz}\) approaches a stationary observer at \(\SI{20}{m.s^{-1}}\). For \(v=\SI{343}{m.s^{-1}}\), the observed frequency is closest to}{%
\item \(\SI{472}{Hz}\)
\item \(\SI{531}{Hz}\)
\item \(\SI{500}{Hz}\)
\item \(\SI{563}{Hz}\)
\item \(\SI{686}{Hz}\)}{%
For an approaching source, \(f'=fv/(v-v_s)=500(343)/(343-20)=\SI{531}{Hz}\).}

\bankq{acoustics-087}{Acoustics: Echoes, Beats, and Doppler Effect}{applied}{D}{acoustics,resonance,pipe}{Calculate a closed-pipe fundamental.}{%
A pipe closed at one end has length \(\SI{0.85}{m}\). Its fundamental frequency is closest to}{%
\item \(\SI{50.4}{Hz}\)
\item \(\SI{80.7}{Hz}\)
\item \(\SI{151}{Hz}\)
\item \(\SI{101}{Hz}\)
\item \(\SI{202}{Hz}\)}{%
For a closed pipe, \(f_1=v/(4L)=343/[4(0.85)]=\SI{101}{Hz}\). Option D carries this value.}

\bankq{acoustics-088}{Acoustics: Echoes, Beats, and Doppler Effect}{applied}{A}{acoustics,intensity,decibel}{Convert an intensity ratio to decibels.}{%
If sound intensity increases by a factor of \(100\), the sound level increases by}{%
\item \(\SI{20}{dB}\)
\item \(\SI{10}{dB}\)
\item \(\SI{40}{dB}\)
\item \(\SI{50}{dB}\)
\item \(\SI{100}{dB}\)}{%
\(\Delta\beta=10\log_{10}(I_2/I_1)=10\log_{10}(100)=\SI{20}{dB}\).}

\bankq{acoustics-089}{Acoustics: Echoes, Beats, and Doppler Effect}{challenge}{C}{acoustics,doppler,reflection}{Analyze the double Doppler shift from a moving reflector.}{%
A stationary detector emits \(\SI{1000}{Hz}\) sound toward a vehicle approaching at \(\SI{30}{m.s^{-1}}\). The reflected frequency received back at the detector is closest to}{%
\item \(\SI{917}{Hz}\)
\item \(\SI{1000}{Hz}\)
\item \(\SI{1192}{Hz}\)
\item \(\SI{1300}{Hz}\)
\item \(\SI{1430}{Hz}\)}{%
The vehicle first acts as moving observer, then as moving source. Thus \(f_r=f(v+u)/(v-u)=1000(343+30)/(343-30)=\SI{1192}{Hz}\).}

\bankq{acoustics-090}{Acoustics: Echoes, Beats, and Doppler Effect}{challenge}{B}{acoustics,echo,falling-object}{Combine free fall with sound travel time.}{%
A stone is dropped into a well and the splash is heard \(\SI{3.00}{s}\) later. Neglecting air resistance, the depth is closest to}{%
\item \(\SI{35.0}{m}\)
\item \(\SI{40.7}{m}\)
\item \(\SI{44.1}{m}\)
\item \(\SI{51.5}{m}\)
\item \(\SI{67.0}{m}\)}{%
Let the fall time be \(t\). Then \(h=\tfrac12gt^2\) and the sound time is \(h/v\), so \(t+h/343=3.00\). Solving gives \(t\approx2.88\) s and \(h\approx\SI{40.7}{m}\).}

% -----------------------------------------------------------------------------
% 91--100: propagation in gases, liquids, and solids
% (2 foundation, 6 applied, 2 challenge)
% -----------------------------------------------------------------------------

\bankq{media-091}{Propagation in Different Media}{foundation}{C}{media,string}{Identify the parameters controlling string-wave speed.}{%
The speed of a transverse wave on an ideal stretched string is}{%
\item \(\sqrt{\mu/\tau}\)
\item \(\tau/\mu\)
\item \(\sqrt{\tau/\mu}\)
\item \(\sqrt{\tau\mu}\)
\item \(\mu/\tau\)}{%
String tension supplies restoring force and linear density supplies inertia, giving \(v=\sqrt{\tau/\mu}\).}

\bankq{media-092}{Propagation in Different Media}{foundation}{A}{media,fluid}{Identify the sound-speed formula in a fluid.}{%
For a fluid of bulk modulus \(B\) and density \(\rho\), the longitudinal sound speed is}{%
\item \(\sqrt{B/\rho}\)
\item \(\sqrt{\rho/B}\)
\item \(B/\rho\)
\item \(\sqrt{B\rho}\)
\item \(1/\sqrt{B\rho}\)}{%
The elastic resistance is measured by \(B\) and inertia by \(\rho\), hence \(v=\sqrt{B/\rho}\).}

\bankq{media-093}{Propagation in Different Media}{applied}{D}{media,string,numerical}{Calculate string-wave speed.}{%
A string has tension \(\SI{100}{N}\) and linear density \(\SI{0.010}{kg.m^{-1}}\). Its wave speed is}{%
\item \(\SI{10}{m.s^{-1}}\)
\item \(\SI{31.6}{m.s^{-1}}\)
\item \(\SI{50}{m.s^{-1}}\)
\item \(\SI{100}{m.s^{-1}}\)
\item \(\SI{1000}{m.s^{-1}}\)}{%
\(v=\sqrt{\tau/\mu}=\sqrt{100/0.010}=\sqrt{10000}=\SI{100}{m.s^{-1}}\).}

\bankq{media-094}{Propagation in Different Media}{applied}{B}{media,liquid,numerical}{Calculate sound speed in a liquid.}{%
Water has \(B=\SI{2.2e9}{Pa}\) and \(\rho=\SI{1000}{kg.m^{-3}}\). Its sound speed is closest to}{%
\item \(\SI{343}{m.s^{-1}}\)
\item \(\SI{1.48e3}{m.s^{-1}}\)
\item \(\SI{2.20e3}{m.s^{-1}}\)
\item \(\SI{4.69e3}{m.s^{-1}}\)
\item \(\SI{2.20e6}{m.s^{-1}}\)}{%
\(v=\sqrt{B/\rho}=\sqrt{2.2\times10^9/1000}=\SI{1.48e3}{m.s^{-1}}\).}

\bankq{media-095}{Propagation in Different Media}{applied}{E}{media,solid,numerical}{Calculate longitudinal wave speed in a rod.}{%
A slender rod has Young's modulus \(Y=\SI{2.0e11}{Pa}\) and density \(\rho=\SI{8.0e3}{kg.m^{-3}}\). Its longitudinal wave speed is closest to}{%
\item \(\SI{500}{m.s^{-1}}\)
\item \(\SI{1.58e3}{m.s^{-1}}\)
\item \(\SI{2.50e3}{m.s^{-1}}\)
\item \(\SI{4.00e3}{m.s^{-1}}\)
\item \(\SI{5.00e3}{m.s^{-1}}\)}{%
\(v=\sqrt{Y/\rho}=\sqrt{2.0\times10^{11}/8.0\times10^3}=\sqrt{2.5\times10^7}=\SI{5.00e3}{m.s^{-1}}\).}

\bankq{media-096}{Propagation in Different Media}{applied}{C}{media,scaling}{Apply stiffness-density scaling.}{%
A material's elastic modulus is quadrupled while its density is unchanged. The wave speed becomes}{%
\item one-quarter as large
\item one-half as large
\item twice as large
\item four times as large
\item unchanged}{%
Since \(v\propto\sqrt{\text{modulus}/\rho}\), multiplying the modulus by four multiplies speed by \(\sqrt4=2\).}

\bankq{media-097}{Propagation in Different Media}{applied}{A}{media,interface}{Identify what remains continuous in transmission.}{%
A periodic wave crosses a stationary boundary into a medium where its speed is lower. Which quantity remains unchanged?}{%
\item frequency
\item wavelength
\item wave number
\item phase velocity
\item amplitude in every case}{%
The boundary is driven at the incident frequency, so transmitted oscillations have the same frequency. Since \(v=f\lambda\), the wavelength changes.}

\bankq{media-098}{Propagation in Different Media}{applied}{D}{media,refraction,waves}{Apply the wave refraction law.}{%
A plane wave crosses from a region of speed \(v_1=\SI{6.0}{m.s^{-1}}\) into one of speed \(v_2=\SI{3.0}{m.s^{-1}}\) at incidence angle \(30^\circ\). The refracted angle is}{%
\item \(7.2^\circ\)
\item \(24.6^\circ\)
\item \(30.0^\circ\)
\item \(14.5^\circ\)
\item \(60.0^\circ\)}{%
Wave refraction gives \(\sin\theta_1/v_1=\sin\theta_2/v_2\). Thus \(\sin\theta_2=(3/6)\sin30^\circ=0.25\), so \(\theta_2=14.5^\circ\), listed as option D.}

\bankq{media-099}{Propagation in Different Media}{challenge}{B}{media,impedance,reflection}{Calculate an amplitude reflection coefficient.}{%
A normally incident acoustic pressure wave crosses from impedance \(Z_1=\SI{2.0e6}{kg.m^{-2}.s^{-1}}\) to \(Z_2=\SI{6.0e6}{kg.m^{-2}.s^{-1}}\). Its pressure-amplitude reflection coefficient is}{%
\item \(-0.50\)
\item \(+0.50\)
\item \(+0.25\)
\item \(+0.75\)
\item \(+2.0\)}{%
For normal incidence, \(r_p=(Z_2-Z_1)/(Z_2+Z_1)=(6-2)/(6+2)=0.50\). The positive sign means no pressure phase reversal.}

\bankq{media-100}{Propagation in Different Media}{challenge}{C}{media,travel-time,multilayer}{Calculate travel time through layered media.}{%
A pulse travels through \(\SI{0.50}{m}\) of water at \(\SI{1500}{m.s^{-1}}\) and then \(\SI{0.80}{m}\) of steel at \(\SI{5000}{m.s^{-1}}\). Ignoring interface delay, the total travel time is}{%
\item \(\SI{0.160}{ms}\)
\item \(\SI{0.333}{ms}\)
\item \(\SI{0.493}{ms}\)
\item \(\SI{0.867}{ms}\)
\item \(\SI{1.30}{ms}\)}{%
Travel times add: \(t=0.50/1500+0.80/5000=3.33\times10^{-4}+1.60\times10^{-4}=4.93\times10^{-4}\) s \(=\SI{0.493}{ms}\).}

\end{quizbank}

\quizselectall
\let\storedquizmarks\quizmarks
\let\storedquizdifficulty\quizdifficulty
\renewcommand{\quizmarks}[1]{}
\renewcommand{\quizdifficulty}[1]{}
\printquizquestions[2]
\let\quizmarks\storedquizmarks
\let\quizdifficulty\storedquizdifficulty

\begin{quizanswerkeycontent}
  \clearpage
  \begin{answerkey}
    \centering
    \renewcommand{\arraystretch}{1.25}
    \begin{tabular}{*{5}{>{\centering\arraybackslash}p{0.17\linewidth}}}
      \toprule
      \textbf{1.} B & \textbf{2.} D & \textbf{3.} A & \textbf{4.} C & \textbf{5.} C \\
      \textbf{6.} E & \textbf{7.} B & \textbf{8.} D & \textbf{9.} A & \textbf{10.} C \\
      \textbf{11.} C & \textbf{12.} B & \textbf{13.} D & \textbf{14.} A & \textbf{15.} C \\
      \textbf{16.} B & \textbf{17.} E & \textbf{18.} C & \textbf{19.} A & \textbf{20.} D \\
      \textbf{21.} B & \textbf{22.} D & \textbf{23.} B & \textbf{24.} D & \textbf{25.} C \\
      \textbf{26.} A & \textbf{27.} D & \textbf{28.} B & \textbf{29.} C & \textbf{30.} E \\
      \textbf{31.} A & \textbf{32.} D & \textbf{33.} B & \textbf{34.} C & \textbf{35.} E \\
      \textbf{36.} C & \textbf{37.} A & \textbf{38.} D & \textbf{39.} B & \textbf{40.} E \\
      \textbf{41.} C & \textbf{42.} A & \textbf{43.} D & \textbf{44.} B & \textbf{45.} E \\
      \textbf{46.} B & \textbf{47.} D & \textbf{48.} A & \textbf{49.} C & \textbf{50.} E \\
      \textbf{51.} B & \textbf{52.} D & \textbf{53.} A & \textbf{54.} C & \textbf{55.} A \\
      \textbf{56.} E & \textbf{57.} B & \textbf{58.} D & \textbf{59.} B & \textbf{60.} C \\
      \textbf{61.} D & \textbf{62.} A & \textbf{63.} E & \textbf{64.} C & \textbf{65.} B \\
      \textbf{66.} D & \textbf{67.} A & \textbf{68.} C & \textbf{69.} D & \textbf{70.} B \\
      \textbf{71.} A & \textbf{72.} C & \textbf{73.} D & \textbf{74.} B & \textbf{75.} E \\
      \textbf{76.} C & \textbf{77.} A & \textbf{78.} D & \textbf{79.} B & \textbf{80.} C \\
      \textbf{81.} B & \textbf{82.} D & \textbf{83.} A & \textbf{84.} C & \textbf{85.} D \\
      \textbf{86.} B & \textbf{87.} D & \textbf{88.} A & \textbf{89.} C & \textbf{90.} B \\
      \textbf{91.} C & \textbf{92.} A & \textbf{93.} D & \textbf{94.} B & \textbf{95.} E \\
      \textbf{96.} C & \textbf{97.} A & \textbf{98.} D & \textbf{99.} B & \textbf{100.} C \\
      \bottomrule
    \end{tabular}
    \medskip

    \textbf{Total marks: 100.}
  \end{answerkey}
\end{quizanswerkeycontent}

% Keep the first solution box on later verso pages below the physical page edge.
% The explicit breaks preserve the class styling while avoiding clipped headers.
\ExplSyntaxOn
\NewDocumentCommand{\printpaginatedquizsolutions}{}
{
  \section{Generated~Worked~Solutions}
  \setcounter{pqselectedquestion}{0}
  \seq_map_inline:Nn \g__pq_selected_ids_seq
  {
    \int_case:nn { \value{pqselectedquestion} }
    {
      {33}{\clearpage\null\par\vspace*{12pt}}
      {55}{\clearpage\null\par\vspace*{12pt}}
      {77}{\clearpage\null\par\vspace*{12pt}}
      {99}{\clearpage\noindent\rule{0pt}{36pt}\par\nobreak}
    }
    \stepcounter{pqselectedquestion}
    \tl_gset:Nn \g__pq_current_id_tl {##1}
    \printsolution{exercise}{##1}
    \tl_gclear:N \g__pq_current_id_tl
  }
}
\ExplSyntaxOff

\begin{quizsolutioncontent}
  \clearpage
  \printpaginatedquizsolutions
\end{quizsolutioncontent}

\begin{quizteachercontent}
  \clearpage
  \printquiztopicreport
  \printquizteacherreport
\end{quizteachercontent}

\end{document}
